\section*{Chapter \ref{ch:b3:water-rock-interaction} --- Water--Rock Interaction}

\begin{solution}{exo:b3:water-rock-interaction:1}
Cations: $2\times48/40.08 + 2\times9.7/24.31 + 23/22.99 + 3.9/39.10 =
\qty{4.29}{meq/L}$; anions: $183/61.02 + 35.5/35.45 + 2\times48/96.06 =
\qty{5.00}{meq/L}$. Error $100\times(4.29 - 5.00)/9.29 = \qty{-7.6}{\percent}$:
far beyond the 1 or \qty{2}{\percent} of a careful analysis of such a water.
A cation is missing or under-reported, or an anion is too high.
\end{solution}

\begin{solution}{exo:b3:water-rock-interaction:2}
$\sqrt I = 0.0707$. \ce{Ca^2+}: $-\log\gamma = 0.5085\times4\times0.0707/(1 +
0.3281\times6\times0.0707) = 0.126$, $\gamma = 0.748$; \ce{HCO3-}: $\gamma =
0.927$; \ce{SO4^2-}: $\gamma = 0.739$.
\end{solution}

\begin{solution}{exo:b3:water-rock-interaction:3}
$\sqrt I = 0.126$; $\gamma_{\ce{Ca}} = 0.622$, $\gamma_{\ce{SO4}} = 0.602$;
$\log\mathrm{IAP} = \log(3.0\times10^{-3}\times0.622\times4.0\times10^{-3}
\times0.602) = -5.35$; $\mathrm{SI} = -5.35 + 4.55 = -0.80$. The water is
undersaturated and can still dissolve gypsum.
\end{solution}

\begin{solution}{exo:b3:water-rock-interaction:4}
$m = 0.012/60.08 = \qty{2.00e-4}{mol/L}$, $\log m = -3.70$. Quartz:
$\mathrm{SI} = -3.70 + 3.98 = +0.28$, slightly supersaturated; amorphous
silica: $-3.70 + 2.71 = -0.99$, undersaturated.
\end{solution}

\begin{solution}{exo:b3:water-rock-interaction:5}
At \qty{25}{\celsius}, $\log K = -8.45 - 6.35 - 1.47 + 10.33 = -5.94$, $m =
(10^{-5.94}\times10^{-2}/4)^{1/3} = \qty{1.42}{mmol/L}$, \qty{57}{mg/L}; at
\qty{10}{\celsius}, $\log K = -5.59$, $m = \qty{1.86}{mmol/L}$,
\qty{75}{mg/L}: a third more in the cold.
\end{solution}

\begin{solution}{exo:b3:water-rock-interaction:6}
$[\ce{CO2}] = 10^{-7.60}\times4.0\times10^{-3}/10^{-6.35} =
\qty{2.25e-4}{mol/L}$; $[\ce{CO3^2-}] = 10^{-10.33}\times4.0\times10^{-3}/
10^{-7.60} = \qty{7.4e-6}{mol/L}$; $\log P_{\ce{CO2}} = -3.648 + 1.47 =
-2.18$, 16 times the air. Total \qty{4.23}{mmol/L}: \ce{CO2}
\qty{5.3}{\percent}, \ce{HCO3-} \qty{94.5}{\percent}, \ce{CO3^2-}
\qty{0.18}{\percent}.
\end{solution}

\begin{solution}{exo:b3:water-rock-interaction:7}
$r = 10^{-12.09} = \qty{8.2e-13}{mol.m^{-2}.s^{-1}}$ (acid, neutral and
basic terms contribute $10^{-12.45}$, $10^{-12.56}$ and $10^{-12.74}$). The
radius shrinks at $rV_m = 8.2\times10^{-13}\times1.0131\times10^{-4} =
\qty{8.3e-17}{m/s}$, \qty{2.6}{nm} a year; the \qty{0.5}{mm} radius lasts
$6.0\times10^{12}$~s, about \num{190000} years.
\end{solution}

\begin{solution}{exo:b3:water-rock-interaction:8}
$\exp[(69.8/0.008314)(1/278.15 - 1/298.15)] = 7.6$; with
\qty{74}{kJ/mol}, 8.6.
\end{solution}

\begin{solution}{exo:b3:water-rock-interaction:9}
Equating $(12.703 - X - 3Y)/3$ and $(7.435 - 2Y)/2$ removes $Y$: $X =
12.703 - 11.153 = 1.55$, a horizontal line. The water has $X = \log(2.0\times
10^{-3}/39.10) + 7.0 = 2.71$ and $Y = \log(0.015/60.08) = -3.60$: above 1.55
and far below the muscovite--K-feldspar line ($X = -3.161 - 3Y = 7.65$), so
it forms muscovite.
\end{solution}

\begin{solution}{exo:b3:water-rock-interaction:10}
$\delta_w^f = -2 - 2 = -4$. Closed: $W/R = 0.5\times(-2 - 6)/(-13 + 4) =
0.44$. Open: $0.5\ln[(-8)/(-9) + 1] = 0.5\ln1.889 = 0.32$. In the open system each
parcel leaves in equilibrium with the rock it last met, carrying away as much
oxygen-18 as it can; in the closed system the whole batch ends at $-4$ and
takes away less per kilogram.
\end{solution}

\begin{solution}{exo:b3:water-rock-interaction:11}
Open: $\log K = -5.59$, $m = (10^{-5.59}\times10^{-1.5}/4)^{1/3} =
\qty{2.73}{mmol/L}$, \qty{109}{mg/L}. Closed: $C_0 = 10^{-1.27 - 1.5} =
\qty{1.70e-3}{mol/L}$ and $K_{\mathrm{cal}}K_1/K_2 = 10^{-4.32}$; the cubic
$4m^3 = 4.79\times10^{-5}(1.70\times10^{-3} - m)$ gives $m =
\qty{1.45}{mmol/L}$, \qty{58}{mg/L}, about half. The carbon dioxide left is
$0.25\times10^{-3}$, \qty{15}{\percent} of $C_0$, so $\log P_{\ce{CO2}}$ falls
from $-1.5$ to $-2.33$.
\end{solution}

\begin{solution}{exo:b3:water-rock-interaction:12}
$1 - 10^{\mathrm{SI}}$: 0.99, 0.90 and 0.21. The power law gives
$(10^4/6)^{-0.61} = 0.011$: a rate about 90 times lower. A field mineral is
old, with its reactive surfaces gone, and bathed in slow pore water close to
equilibrium; both factors multiply, and the field rate falls orders of
magnitude below the laboratory rate.
\end{solution}

\begin{solution}{pb:b3:water-rock-interaction:1}
\textbf{1.} \ce{Ca} 2.795, \ce{Mg} 0.165, \ce{Na} 0.305, \ce{K} 0.031,
\ce{HCO3} 5.572, \ce{Cl} 0.310, \ce{SO4} 0.146, \ce{NO3} \qty{0.145}{mmol/L}.

\textbf{2.} Cations \qty{6.25}{meq/L}, anions \qty{6.32}{meq/L}: error
\qty{-0.5}{\percent}, acceptable.

\textbf{3.} $I = \tfrac12(4\times2.795 + 4\times0.165 + 0.305 + 0.031 + 5.572
+ 0.310 + 4\times0.146 + 0.145)\times10^{-3} = 0.0094$.

\textbf{4.} $\sqrt I = 0.0969$: $\gamma_{\ce{Ca^2+}} = 0.683$,
$\gamma_{\ce{HCO3-}} = 0.904$, $\gamma_{\ce{CO3^2-}} = 0.676$.

\textbf{5.} Yes: $I$ is a tenth of the limit of about 0.1.

\textbf{6.} $a_{\ce{HCO3-}} = 5.572\times10^{-3}\times0.904 = 5.04\times
10^{-3}$ ($\log = -2.298$).

\textbf{7.} $a_{\ce{CO2}} = 10^{-7.20}\times5.04\times10^{-3}/10^{-6.46} =
9.17\times10^{-4}$ ($\log = -3.038$).

\textbf{8.} $\log P_{\ce{CO2}} = -3.038 + 1.27 = -1.77$, \qty{0.017}{atm}:
$10^{1.60} = 40$ times the air, within the ground-water range. The carbon
dioxide comes from the soil air of the recharge area, from roots and decay.

\textbf{9.} $a_{\ce{CO3^2-}} = 10^{-10.49}\times5.04\times10^{-3}/10^{-7.20}
= 2.58\times10^{-6}$, $[\ce{CO3^2-}] = \qty{3.8e-6}{mol/L}$. Total
\qty{6.49}{mmol/L}: \ce{CO2} \qty{14.1}{\percent}, \ce{HCO3-}
\qty{85.8}{\percent}, \ce{CO3^2-} \qty{0.06}{\percent}.

\textbf{10.} Alkalinity $5.572 + 2\times0.0038 = \qty{5.580}{meq/L}$:
\qty{99.9}{\percent} bicarbonate.

\textbf{11.} $a_{\ce{Ca^2+}} = 2.795\times10^{-3}\times0.683 = 1.909\times
10^{-3}$; $\log\mathrm{IAP} = \log(1.909\times10^{-3}\times2.58\times10^{-6})
= -8.307$; $\mathrm{SI} = -8.307 + 8.35 = +0.04$.

\textbf{12.} Within $\pm0.1$, the uncertainty of the index: the water is in
equilibrium with calcite. It has stayed long in the rock: a slow, diffuse
spring.

\textbf{13.} $m = [10^{-5.59}\times0.0171/(4\times0.683\times0.904^2)]^{1/3}
= \qty{2.70}{mmol/L}$, \qty{108}{mg/L}, against \qty{2.79}{mmol/L} measured:
\qty{3}{\percent} more, as an index of $+0.04$ means.

\textbf{14.} $a_{\ce{CO2}} = 10^{-1.27 - 3.37} = 2.29\times10^{-5}$;
$a_{\ce{H+}} = 10^{-6.46}\times2.29\times10^{-5}/5.04\times10^{-3}$, pH 8.80;
$a_{\ce{CO3^2-}} = 1.03\times10^{-4}$; $\mathrm{SI} = \log(1.909\times
10^{-3}\times1.03\times10^{-4}) + 8.35 = +1.65$ (1.61 with the full
alkalinity balance of the figure).

\textbf{15.} Iterating $m = [10^{-5.59 - 3.37}/(4\gamma_{\ce{Ca}}
\gamma_{\ce{HCO3}}^2)]^{1/3}$ with $I = 3m$: $m = \qty{0.72}{mmol/L}$
(\qty{29}{mg/L}), $I = 0.0022$. Calcite precipitated: $2.795 - 0.718 =
\qty{2.08}{mmol/L}$, about \qty{210}{mg} per litre.

\textbf{16.} $0.208\times2.0 = \qty{0.42}{g/s}$, \qty{13}{t} a year,
$13.1/2.71 = \qty{4.8}{m^3}$; over \qty{400}{m^2}, \qty{12}{mm} a year,
\qty{0.03}{mm} a day: well below the cold-spring tufas quoted in the chapter.

\textbf{17.} $X = \log(1.2\times10^{-3}/39.10) + 5.8 = 1.29$; $Y =
\log(9.0\times10^{-3}/60.08) = -3.82$.

\textbf{18.} $X < 1.55$ and $Y > -4.39$: kaolinite. Quartz: $\mathrm{SI} =
-3.82 + 3.98 = +0.16$.

\textbf{19.} Below $\log a_{\ce{H4SiO4}} = -4.39$, \qty{2.4}{mg/L} of
silica.

\textbf{20.} $X = 1.55$ needs pH $= 1.55 + 4.51 = 6.06$.

\textbf{21.} Acid $-10.06 - 0.500\times5.8 = -12.96$; neutral $-12.41$; basic
$-21.20 + 0.823\times5.8 = -16.43$. Sum: $\log r = -12.30$, the neutral term
dominant.

\textbf{22.} Supply-limited: thick soil and gentle slopes mean slow erosion;
the fresh feldspar that reaches the weathering front is fully leached, and
the flux follows the erosion rate rather than the climate.

\textbf{23.} $\exp[(74/0.008314)(1/283.15 - 1/298.15)] = 4.9$.

\textbf{24.} The calcite saturation index of the spring water at its outlet
is $\mathrm{SI} = \textbf{+0.04}$.
\end{solution}
